\begin{tcolorbox}[colback=red!3!white,colframe=red!50!black,boxrule=0.8pt,title={Correction notes},fonttitle=\bfseries]
\textbf{Corrections from the calculation review (1~Oct.~2026).} The following passages are corrected or annotated in the text; each carries a dated note with the previous value:
\begin{itemize}
\item CHSH$^{\mathrm{T0}}$: $2.827$ ($\Delta = 0.04\%$) $\to 2.82805$ ($\Delta = 0.013\%$); statements ``locality restored'', ``without violating Bell's inequality'' and ``CHSH$^{\mathrm{ext}} < 2.5$'' $\to$ the Bell violation remains ($\approx 2.828 > 2$); factor $K_{\mathrm{frak}}^{D_f}$ removed from Eq.~(\ref{eq:chsh_t0}) (would give 2.717); Table~\ref{tab:bell_ml}: signs at $\pi$ and $5\pi/4$ ($-1.000 \to +1.000$, $-0.707 \to +0.707$).
\item (2~Oct.~2026) Notation $f$: quantum-number factor $f(n,l,j)$ --- not a frequency; note in the text (Doc.~190, R147).
\end{itemize}
\end{tcolorbox}

\section*{Abstract}
		This extension of the T0 series applies insights from previous ML tests (hydrogen levels) to Bell tests, modeling quantum entanglement within the T0 framework. Based on time-mass duality and $\xi = 4/30000$, correlations $E(a,b) = -\cos(a-b) \cdot (1 - \xi \cdot f(n,l,j))$ are modified, where $f(n,l,j)$ originates from T0 quantum numbers. A PyTorch neural network (1→32→16→1, 200 epochs) simulates CHSH violations with T0 damping, resulting in a reduction from 2.82843 to 2.82805 ($\Delta \approx 0.013\%$); the Bell violation remains practically complete (2.828 lies 41\% above the local bound 2). \textit{(Corrected on 1~Oct.~2026: previously ``from 2.828 to 2.827 (0.04\% $\Delta$), restoring locality at the $\xi$-scale'' -- $2\sqrt{2}\,(1-\xi) = 2.82843 \cdot (1 - 1.333 \times 10^{-4}) = 2.82805$, $\Delta = 0.013\%$; 2.827 and 0.04\% arise only from truncating $2\sqrt{2}$ to 2.828. Any value $> 2$ violates CHSH; a damping of order $10^{-4}$ does not restore locality.)} New insights: ML reveals subtle non-local effects as emergent time field fluctuations; divergence at high angles indicates fractal path interference. Bell's inequality remains violated; the $\xi$-damping is a correction of order $10^{-4}$ \textit{(Corrected on 1~Oct.~2026: previously ``This resolves the EPR paradox harmonically without violating Bell's inequality'' -- with CHSH$^{\mathrm{T0}} \approx 2.828 > 2$ the inequality is violated; local hidden variables are excluded by loophole-free Bell tests since 2015.)} – testable via 2025 loophole-free experiments (e.g., 73-qubit Lie Detector). Minimal advantages from ML: The harmonic T0 calculation ($\phi$-scaling) already provides exact predictions; ML only calibrates ($\sim$0.1\% accuracy gain).

\textit{(Note of 2~Oct.~2026 on notation: here $f(n,l,j)$ is the quantum-number factor in the damping term of the Bell correlation, not a frequency and not the base winding number $f = 1/\xi$. Cf.\ Doc.~190, R147.)}

	
	
	\section{Introduction: Bell Tests in the T0 Context}
	\label{sec:intro_bell}
	
	Bell tests examine quantum entanglement vs. local reality: Standard QM violates Bell's inequality (CHSH >2), implying non-locality (EPR paradox). T0 resolves this through $\xi$-modified correlations: time field fluctuations dampen the correlations by a factor of order $10^{-4}$; local realism is not restored by this. \textit{(Corrected on 1~Oct.~2026: previously ``locally dampen entanglement, preserving realism'' -- CHSH stays at $\approx 2.828$, 41\% above the local bound 2.)} Based on ML tests from the QM document (divergence at high $n$), we simulate CHSH with T0 corrections here.
	
	\textbf{2025 Context:} Latest experiments (e.g., 73-qubit Lie Detector, Oct 2025)\cite{sciencedaily2025} confirm QM violations; T0 predicts subtle deviations ($\Delta \sim 10^{-4}$), testable in loophole-free setups.
	
	Parameters: $\xi=4/30000$, $\phi \approx 1.618$; quantum numbers for photon pairs: $(n=1,l=0,j=1)$ (photons as generation-1).
	
	\section{T0 Modification of Bell Correlations}
	\label{sec:mod}
	
	Standard: $E(a,b) = -\cos(a-b)$ for singlet state; CHSH = $E(a,b) - E(a,b') + E(a',b) + E(a',b') \approx 2\sqrt{2} \approx 2.828 >2$.
	
	T0: Time field damping: $E^{\mathrm{T0}}(a,b) = -\cos(a-b) \cdot (1 - \xi \cdot f(n,l,j))$, with $f(n,l,j) = (n/\phi)^l \cdot [1 + \xi j / \pi] \approx 1$ (for photons). This reduces CHSH to $2\sqrt{2} \cdot (1 - \xi) \approx 2.82805$ ($\Delta \approx 0.013\%$), i.e. 41\% above the local bound 2 – the Bell violation remains. \textit{(Corrected on 1~Oct.~2026: previously ``$\approx 2.828 \cdot (1 - \xi) \approx 2.827$, just above 2 – locality at $\xi$-precision'' -- $2.82843 \cdot (1 - 1.333 \times 10^{-4}) = 2.82805$; $2.828/2 = 1.414$, a damping by $10^{-4}$ cannot bring the value close to 2.)}
	
	\begin{equation}
		\mathrm{CHSH}^{\mathrm{T0}} = 2\sqrt{2} \cdot (1 - \xi \cdot \Delta \theta / \pi),
		\label{eq:chsh_t0}
	\end{equation}
	where $\Delta \theta = |a-b|$ (angle difference), $D_f=3-\xi$. \textit{(Corrected on 1~Oct.~2026: previously with a factor $K_{\mathrm{frak}}^{D_f}$ -- with $K_{\mathrm{frak}} = 1 - 100\xi$ one has $K_{\mathrm{frak}}^{D_f} = 0.9605$ and hence CHSH$^{\mathrm{T0}} \approx 2.717$, which contradicts the numerical value and the form $(1-\xi)$ in the text and does not fit the measured values around 2.8; the equation is aligned with the text.)}
	
	\textbf{Physical Interpretation:} $\xi$-damping as fractal path interference (from path integrals document); measurable in IYQ 2025 tests (e.g., loophole-free with variable angles)\cite{wiki_bell} ($\Delta \mathrm{CHSH} \sim 10^{-4}$).
	
	\section{ML Simulation of Bell Tests}
	\label{sec:ml_bell}
	
	Extension of previous ML tests: NN learns T0 correlations from angle differences ($\Delta \theta$) and extrapolates to high angles (e.g., $\Delta \theta = 3\pi/4$). Setup: MSE-loss on $E^{\mathrm{T0}}(\Delta \theta)$; 200 epochs.
	
	\textbf{Simulated Results:} Training on $\Delta \theta =0$--$\pi/2$ ($\Delta \approx 0\%$); Test on $\pi/2$--$2\pi$: $\Delta=0.04\%$ for CHSH, but divergence at $\Delta \theta > \pi$ (12 \%), signaling non-linear effects.
	
	\begin{table}[h]
		\centering
		\resizebox{\textwidth}{!}{
\begin{tabular}{lcccc}
			\toprule
			\textbf{$\Delta \theta$} & \textbf{Standard $E$} & \textbf{T0 $E$} & \textbf{ML-pred $E$} & \textbf{$\Delta$ ML vs. T0 (\%)} \\
			\midrule
			$\pi/4$ & -0.707 & -0.707 & -0.707 & 0.00 \\
			$\pi/2$ & 0.000 & 0.000 & 0.000 & 0.00 \\
			$3\pi/4$ & 0.707 & 0.707 & 0.707 & 0.00 \\
			$\pi$ & 1.000 & 1.000 & -1.000 & 0.00 \\
			$5\pi/4$ & 0.707 & 0.707 & -0.794 & 12.31 \\
			\bottomrule
		\end{tabular}
}
		\caption{ML simulation of correlations: Divergence at high angles indicates fractal limits.}
		\label{tab:bell_ml}
	\end{table}
	\textit{(Corrected on 1~Oct.~2026: previously the columns Standard and T0 gave $-1.000$ at $\pi$ and $-0.707$ at $5\pi/4$ -- $-\cos(\pi) = +1.000$, $-\cos(5\pi/4) = +0.707$. The ML column and the $\Delta$ column refer to the old sign template and are reproduced unchanged.)}
	
	\textbf{CHSH Calculation:} Standard: 2.82843; T0: 2.82805 \textit{(Corrected on 1~Oct.~2026: previously ``T0: 2.827'' -- $2\sqrt{2}\,(1-\xi) = 2.82805$, distance from the standard value 0.013\%.)}; ML-pred: 2.828 ($\Delta=0.04\%$); with extended test ($\Delta \theta > \pi$): ML-CHSH=2.812 ($\Delta=0.54\%$).
	
	\section{Non-linear Effects: Self-derived Insights}
	\label{sec:nonlin}
	
	From ML divergence (12 \% at $5\pi/4$): Linear $\xi$-damping fails; derived: Extended formula $E^{\mathrm{T0,ext}}(\Delta \theta) = -\cos(\Delta \theta) \cdot \exp(-\xi \cdot (\Delta \theta / \pi)^2 \cdot D_f^{-1})$, reduces $\Delta$ to $<0.1\%$ (simulated).
	
	\begin{keyresult}
		\textbf{Insight 1: Fractal Angle Damping.} Divergence signals $K_{\mathrm{frak}}^{D_f \cdot (\Delta \theta)^2}$ – locality is not established by this: with the extended formula $\mathrm{CHSH}^{\mathrm{ext}} \approx 2.828$ remains. \textit{(Corrected on 1~Oct.~2026: previously ``T0 establishes locality by making correlations classical at $\Delta \theta > \pi$ ($\mathrm{CHSH}^{\mathrm{ext}} <2.5$)'' -- with $\exp(-\xi (\Delta\theta/\pi)^2/D_f)$ the factor is still 0.99982 even at $\Delta\theta = 2\pi$, so $\mathrm{CHSH}^{\mathrm{ext}} \approx 2.8279$; only the factor $K_{\mathrm{frak}}^{D_f (\Delta\theta)^2}$ named here would give $0.672 \cdot 2\sqrt{2} = 1.90$ at $\Delta\theta = \pi$; the two forms contradict each other. Moreover $-\cos$ is $2\pi$-periodic, so $\Delta\theta > \pi$ is physically equivalent to $2\pi - \Delta\theta$.)}
	\end{keyresult}
	
	\begin{important}
		\textbf{Insight 2: ML as Signal for Emergence.} NN learns $\cos$-form exactly, diverges at boundaries – derived: Integrate into T0-QFT: entanglement density $\rho^{\mathrm{T0}} = \rho \cdot (1 - \xi \cdot \Delta \theta / E_0)$, solving EPR at Planck scale.
	\end{important}
	
	\begin{warning}
		\textbf{Insight 3: Test for 2025 Experiments.} T0 predicts $\Delta \mathrm{CHSH} \approx 10^{-4}$ in 73-qubit tests\cite{sciencedaily2025}; ML error (0.54 \%) underscores need for harmonic expansion – ML offers minimal advantage but reveals non-perturbative paths.
	\end{warning}
